\centerline{\bf \Huge ОТА. Делимость II}

\begin{flushright}
        \scriptsize{Выживает тот, у кого есть воля к жизни.\\ Мисато Кацураги}
\end{flushright}

\problem{1}
Перемножили несколько натуральных чисел и получили 224, причём самое маленькое число было ровно вдвое меньше самого большого. Сколько чисел перемножили?

\textbf{Решение.}
$224=2^5 \cdot 7$. Рассмотрим два числа, указанные в условии: самое маленькое и самое большое. Если одно из них делится на 7, то и другое должно образом, оно равно $2^3=8$, а самое маленькое - это $2^2=4$. Остался единственный множитель, равный семи, поэтому искомых чисел три: 4,7 и 8 .

\problem{2}
Существует ли целое число, произведение цифр которого равно  а) 1980?  б) 1990?  в) 2000?

\textbf{Решение.}
Среди делителей числа $1980$ --- двузначное простое число $11$, а среди делителей числа $1990$ --- трёхзначное простое число $199$, поэтому ни $1980$, ни $1990$ в такие произведения разложить нельзя. Число $2000$ можно разными способами разложить на однозначные множители, например так:  $2000 = 2\cdot4\cdot5^3$.  Следовательно, таких чисел, как требуется в условии, достаточно много. Вот, например, два из них – $555422$ и $25855$.

\problem{3}
По кругу расставили 2017 чисел. Может ли быть так, что отношение любых двух последовательных чисел - простое число?

\textbf{Решение:}
Допустим, что нашлись числа $a_1, a_2, \ldots, a_{2017}$, которые можно расставить требуемым образом. Пусть число $a_k(k=1,2, \ldots, 2017)$ представляется в виде произведения $n_k$ простых сомножителей (не обязательно различных). Так как каждые два соседние числа отличаются друг от друга одним простым множителем, то для каждого $k=1,2, \ldots, 2016$ числа $n_k$ и $n_{k+1}$ отличаются на единицу, то есть имеют разную чётность. Значит, числа $n_1, n_3, \ldots, n_{2017}$ должны быть одной чётности. С другой стороны, числа $a_{2017}$ и $a_1$ также соседние, поэтому $n_{2017}$ и $n_1$ должны иметь разную чётность. Противоречие.

\textit{Для ассистентов:} подскажите им следить за количеством простых чисел в разложении, а именно за чётностью этого количества. Попросить узнать, что происходит с переходом на другое число по кругу.

\problem{4}
Натуральное число назовём удивительным, если самый большой его собственный делитель (т.е. делитель, не равный 1 и самому числу) на 1 больше самого маленького собственного делителя. Найдите все удивительные числа.

\textbf{Решение.} Пусть $p$ самый маленький простой делитель числа $n$. Тогда $\dfrac{n}{p}$ самый большой. По условию $p+1 = \dfrac{n}{p}$. То есть $n = p(p+1)$, где $p$ простое число. 

\textit{Для ассистентов:} попросите доказать, что самый маленький собственный делитель числа это простое число, а самый большой это само число делённое на это простое. Дальше попросите записать уравнение.

\problem{5}
Опишите все натуральные числа, у которых нечётное количество делителей.

\textbf{Решение.} Заметим, что все делители числа $n$ разбиваются на пары по типу $(d, \dfrac{n}{d})$. Но у делителя не будет пары тогда и только тогда когда он сам себе пара, то есть $d = \dfrac{n}{d}$, то есть $n = d^2$. Стало быть это точные квадраты.

\textit{Для ассистентов:} подскажите ученику идею разбиения делителей на пары.

\problem{6}
На доске записаны двузначные числа. Каждое число составное, но любые два числа взаимно просты. Какое наибольшее количество чисел может быть записано?

\textbf{Решение.}
Оценка. Так как любые два записанных числа взаимно просты, то каждое из простых чисел $2,3,5$ и $7$ может войти в разложение на множители не более, чем одного из них. Если на доске пять или более чисел, то все простые множители в разложении какого-то из них должны быть не меньше чем 11 . Но это составное число, значит, оно не меньше чем 121 . Это противоречит условию. Следовательно, на доске записано не более четырёх чисел.

Пример. 25, 26, 33, 49.